\scanpage{60}
% archive sheet number: 59
$D$, \deleted{on sait que} la question étant locale \inserted{sur $Y$}, on peut supposer que $\underline{L}\simeq\underline{O}_X(k)$ par la \unsure{remarque}, d'autre part on \unsure{voit|sait} aussitôt, puisque $\underline{L}\mid f^{-1}(y)\simeq\underline{O}_{f^{-1}(y)}(k)$, que $k=n$. Donc $D$ est le diviseur d'une section de $\underline{O}_X(\unsure{n|k})$.

Ainsi, on a \deleted{\unsure{prouvé}} trouvé un isomorphisme \inserted{\unsure{ensembliste}}
\[
\underline{\unsure{\Phi|\varphi}}^{(n)}(P(\underline{E}))(Y)\xrightarrow{\sim}\unsure{\mathscr{D}|D}^{(n)}(X/Y).
\]
Cet isomorphisme est fonctoriel, compatible aux images inverses.

D'autre part, la notion de ``diviseur \inserted{loc. principal positif} transversal'' est susceptible de descente [car c'est la notion de : fibré inversible, muni d'une section régulière qui est non div. de $0$ et non diviseur de $0$ sur chaque fibre]. De plus, la notion de degré d'un diviseur sur un schéma algébrique de \unsure{Sévéri|Séveri}-Brauer a un sens \deleted{[\unsure{on étend la base} \inserted{\unsure{alg. clos}}]}, invariant par extension du corps de base, d'où la notion des \uline{diviseurs \inserted{positifs} transversaux de degré $n$} sur un \unsure{sch. Brauer}. Cette notion est susceptible de descente. D'où le

\uline{\unsure{Théorème}} Les sections de $\Phi^{(n)}(X)$ ($X$ un \unsure{Sév. Br.}) sont les diviseurs \inserted{positifs} transversaux de degré $n$ de $X/Y$.
